\section{Classification of rank-one harmonic matrix pencils}
\label{sec:classification}

\subsection{The rational-weight model, including several lattices}
Let $K\subseteq\C$ be a characteristic-zero field containing the
parameters below. Consider a proper rational function
\begin{equation}\label{eq:rational-weight}
 R(x)=\sum_{\nu=1}^{L}\sum_{j=0}^{J_\nu}
             \sum_{h=1}^{m_\nu}\frac{c_{\nu,j,h}}{(x+a_\nu+j)^h},
 \qquad \Rea a_\nu>0.
\end{equation}
Representatives from distinct lattices are chosen inequivalent modulo
$\mathbb Z$, so that the associated harmonic families are not artificially
duplicated. Zero coefficients and gaps in a lattice are allowed.
For $q\ge1$ put
\begin{equation}\label{eq:shifted-harmonic}
 H_N^{(q)}(a_\nu)=\sum_{n=0}^N(n+a_\nu)^{-q}.
\end{equation}
The finite moment matrix is
\begin{equation}\label{eq:rational-moment-matrix}
 M_N=\left(\sum_{n=0}^Nn^{i+j}R(n)\right)_{0\le i,j<r}.
\end{equation}

\begin{proposition}[Explicit harmonic coefficient matrices]\label{prop:pencil}
There is a matrix $D_N$ with entries rational in $N$ over $K$ such that
\begin{equation}\label{eq:pencil-decomposition}
 M_N=D_N+\sum_{\nu=1}^L\sum_{q=1}^{m_\nu}
                H_N^{(q)}(a_\nu)\,C_{\nu,q}.
\end{equation}
Each $C_{\nu,q}$ is a fixed symmetric Hankel matrix. Writing
$(C_{\nu,q})_{ij}=c_{\nu,q}[i+j]$, its entries are
\begin{equation}\label{eq:Centries}
 c_{\nu,q}[k]=\sum_{h=q}^{m_\nu}\binom{k}{h-q}
       \sum_{j=0}^{J_\nu}c_{\nu,j,h}(-a_\nu-j)^{k-h+q},
       \quad 0\le k\le d.
\end{equation}
A summand is omitted unless $0\le h-q\le k$.
\end{proposition}
\begin{proof}
Expand $x^k=((x+a_\nu+j)-(a_\nu+j))^k$. The coefficient of
$(x+a_\nu+j)^{-q}$ in $x^kR(x)$ is exactly the corresponding summand
in \eqref{eq:Centries}. Summing over $n=0,\ldots,N$ and shifting the
finite harmonic sum changes it by
\[
 -\sum_{\ell=0}^{j-1}(\ell+a_\nu)^{-q}
 +\sum_{\ell=N+1}^{N+j}(\ell+a_\nu)^{-q},
\]
which is rational in $N$ over $K$. The polynomial part of $x^kR(x)$
sums to a polynomial in $N$, for instance by the finite Bernoulli-polynomial
summation formula. All remaining terms therefore belong to $D_N$.
\end{proof}

\begin{definition}
The \emph{harmonic matrix pencil} associated with $R$ and $r$ is
\begin{equation}\label{eq:harmonic-pencil}
 \mathscr C(\bm X)=\sum_{\nu,q}X_{\nu,q}C_{\nu,q},
\end{equation}
where the variables $X_{\nu,q}$ are algebraically independent. Its rank
is the generic matrix rank over $K(\bm X)$.
\end{definition}
This definition is intrinsic after the distinct shift-orbits have been
fixed. It does not posit algebraic independence of the actual functions
$H_N^{(q)}(a_\nu)$ or of their limits.

\subsection{The complete rank-one test}
For $\tau\in K$, write
\[
 v(\tau)=(1,\tau,\ldots,\tau^{r-1})^{\mathsf T},\qquad
 \ee=(0,\ldots,0,1)^{\mathsf T}.
\]

\begin{theorem}[All rank-one harmonic pencils]\label{thm:classification}
The nonzero pencil \eqref{eq:harmonic-pencil} has rank one if and only
if exactly one of the following forms is possible:
\begin{align}
 C_{\nu,q}&=\lambda_{\nu,q}v(\tau)v(\tau)^{\mathsf T}
       &&\text{for a common finite }\tau,\label{eq:geometric-branch}\\
 C_{\nu,q}&=\lambda_{\nu,q}\ee\ee^{\mathsf T}
       &&\text{for all }\nu,q.\label{eq:corner-branch}
\end{align}
In either case some $\lambda_{\nu,q}$ is nonzero. Equivalently, the
finite residue sums \eqref{eq:Centries} obey either
\[
 c_{\nu,q}[k]=\lambda_{\nu,q}\tau^k\quad(0\le k\le d)
\]
with one common $\tau$, or
\[
 c_{\nu,q}[k]=\lambda_{\nu,q}\mathbf1_{k=d}\quad(0\le k\le d).
\]
These are necessary and sufficient algebraic tests for the original
rational weight, including repeated poles and multiple lattices.
\end{theorem}
\begin{proof}
If a symmetric pencil has rank at most one, every coefficient matrix
has rank at most one. Two nonzero symmetric rank-one matrices can be
written $\alpha uu^{\mathsf T}$ and $\beta vv^{\mathsf T}$ over a
field extension. If $u,v$ are not proportional, some pair of coordinates
has $u_iv_j-u_jv_i\ne0$. The corresponding principal minor of
$\alpha uu^{\mathsf T}+t\beta vv^{\mathsf T}$ is
\[
 t\alpha\beta(u_iv_j-u_jv_i)^2\ne0.
\]
This contradicts rank at most one. All coefficient matrices consequently
have one common rank-one direction.

It remains to classify a nonzero rank-one symmetric Hankel matrix.
Write it as $\lambda uu^{\mathsf T}$. If $u_0\ne0$, normalize
$u_0=1$. Equality of the entries $(0,j+1)$ and $(1,j)$ gives
$u_{j+1}=u_1u_j$, hence $u_j=\tau^j$. If $u_0=0$, let $k$ be
the first nonzero index. Were $k<r-1$, the Hankel entries $(k,k)$
and $(k-1,k+1)$ would respectively be nonzero and zero. Thus
$k=r-1$, giving the corner direction. The finite parameter $\tau$
belongs to $K$: it is the ratio of entries $(0,1)$ and $(0,0)$ of
any coefficient matrix whose $(0,0)$ entry is nonzero. The converse
is immediate from either displayed form.
\end{proof}

The test is constructive. If some $C_{\nu,q}(0,0)\ne0$, compute
$\tau=C_{\nu,q}(0,1)/C_{\nu,q}(0,0)$ and verify
\eqref{eq:geometric-branch} in all entries. If no such coefficient
exists, only \eqref{eq:corner-branch} can hold. A zero pencil is
reported separately, not mislabeled as rank one.

\subsection{The critical-decay theorem: repeated finite differences}
Let the residue polynomial for pole order $q$ on lattice $\nu$ be
\begin{equation}\label{eq:Pq}
 P_{\nu,q}(z)=\sum_{j=0}^{J_\nu}c_{\nu,j,q}z^j.
\end{equation}

\begin{lemma}[Corner support and finite-difference divisibility]
\label{lem:corner-divisibility}
The following statements are equivalent, with no critical-decay assumption:
\begin{enumerate}
\item Every $C_{\nu,q}$ is supported in its bottom-right entry.
\item For every $\nu,q$ and $0\le k<d$,
      $\sum_jc_{\nu,j,q}(-a_\nu-j)^k=0$.
\item Every $P_{\nu,q}(z)$ is divisible by $(z-1)^d$ in $K[z]$.
\item There exists a unique proper rational function $F$ on the same
      finite collection of translated lattices such that
      $R=\Delta^dF/d!$.
\end{enumerate}
The harmonic pencil is zero exactly when every $P_{\nu,q}$ is
divisible by $(z-1)^{d+1}$, equivalently when $R$ is a
$(d+1)$st finite difference of a proper rational function.
\end{lemma}
\begin{proof}
Assume corner support. At the highest pole order $q=m_\nu$,
\eqref{eq:Centries} states directly that the moments of the residues
of order $q$ vanish through degree $d-1$. Descend in $q$.
Every already treated term with $h>q$ involves a residue moment of
degree $k-h+q<k<d$ and hence vanishes. The remaining term is exactly
the required order-$q$ moment. This proves the implication from the
first to the second statement. Substitution in \eqref{eq:Centries}
proves the converse.

Shifting by $a_\nu$ does not change the span of polynomials of degree
less than $d$. Ordinary powers of $j$ and the falling factorials
$\fall{j}{k}$ span that same space. Thus the moment conditions are
equivalent to $P_{\nu,q}^{(k)}(1)=0$ for $k<d$, which is precisely
the asserted divisibility.

For a proper rational function $F$ with residue polynomials
$Q_{\nu,q}$, replacing $x$ by $x+1$ multiplies each residue polynomial
by $z$. Therefore $\Delta^dF/d!$ has residue polynomials
$(z-1)^dQ_{\nu,q}/d!$. Dividing $P_{\nu,q}$ explicitly constructs
$F$ and proves the fourth statement. Multiplication by $(z-1)^d$
is injective on polynomial residue data, so this proper rational $F$
is unique. A polynomial kernel is excluded by properness.

For a zero pencil repeat the same descending argument through degree
$d$, rather than $d-1$. This gives divisibility by $(z-1)^{d+1}$
and its stated finite-difference interpretation.
\end{proof}

\begin{theorem}[Critical rational weights: an exact if-and-only-if]
\label{thm:critical}
Suppose $R$ has the form \eqref{eq:rational-weight} and
\begin{equation}\label{eq:critical-decay}
 R(x)=\kappa x^{-(d+1)}+O(x^{-(d+2)}),\qquad\kappa\ne0,
 \qquad d=2r-2.
\end{equation}
Then its harmonic pencil has rank one if and only if
\begin{equation}\label{eq:critical-difference}
 \boxed{\displaystyle R(x)=\frac1{(2r-2)!}\Delta^{2r-2}F(x)}
\end{equation}
for a proper rational function $F$ on the same translated lattices.
Equivalently, every residue polynomial \eqref{eq:Pq} is divisible
by $(z-1)^{2r-2}$. In that event every harmonic coefficient lies in
the single bottom-right entry of the moment matrix.
\end{theorem}
\begin{proof}
The coefficient of $x^{-1}$ in the expansion of $x^kR(x)$ at infinity
is the sum of the coefficients of all its simple partial-fraction
terms. By \eqref{eq:critical-decay}, it is zero for $k<d$ and
$\kappa$ for $k=d$. Proposition~\ref{prop:pencil} therefore gives
\begin{equation}\label{eq:sumC1}
 \sum_{\nu=1}^L C_{\nu,1}=\kappa\ee\ee^{\mathsf T}.
\end{equation}
If the harmonic pencil has rank one, the common-image argument in
Theorem~\ref{thm:classification} forces all its coefficient matrices
to be multiples of this nonzero corner matrix. Apply
Lemma~\ref{lem:corner-divisibility}. Conversely, the finite-difference
condition gives corner support. Equation~\eqref{eq:sumC1} guarantees
that the pencil is nonzero, hence has rank exactly one.
\end{proof}

This answers the source's repeated-pole rank-one question at its
critical decay and supplies a complete finite residue test even away
from that decay. The extension to several residue classes does not
require a new determinant argument: equation~\eqref{eq:sumC1} sums
over those classes, while the finite-difference divisibility is checked
separately on each shift-orbit.

If
\[
 F(x)=\sum_{\nu,j,q}\frac{b_{\nu,j,q}}{(x+a_\nu+j)^q},
 \qquad R=\Delta^dF/d!,
\]
then the corner coefficient is especially simple:
\begin{equation}\label{eq:lambda-seed}
 (C_{\nu,q})_{r-1,r-1}=\sum_j b_{\nu,j,q}.
\end{equation}
Indeed, the $d$th residue moment of a normalized $d$th difference
returns the leading coefficient of that polynomial moment; terms
of higher pole order involve lower residue moments and vanish.
For $F(x)=(x+a)^{-m}$ this coefficient is one at order $m$ and zero
at every other order. This recovers exactly the moment cancellation
of Section~\ref{sec:spectral} at integer spectral orders.

\subsection{Why the nonzero critical coefficient matters}
\label{sec:counterexample}
For $r=2$, consider
\begin{align}
 R(x)&=\frac{2}{(x+1)^2(x+2)(x+3)}\label{eq:counterexample}\\
 &=\frac1{(x+1)^2}-\frac{3}{2(x+1)}
       +\frac2{x+2}-\frac1{2(x+3)}.\notag
\end{align}
It decays as $2x^{-4}$, faster than the critical $x^{-3}$.
Its harmonic coefficient matrices are
\[
 C_1=0,\qquad
 C_2=\begin{pmatrix}1&-1\\-1&1\end{pmatrix}
       =v(-1)v(-1)^{\mathsf T}.
\]
The harmonic pencil has rank one, but $P_2(z)=1$ is not divisible
by $(z-1)^2$. Thus $R$ is not a second difference of a proper
rational function. This is the geometric branch in
Theorem~\ref{thm:classification}, not a counterexample to
Theorem~\ref{thm:critical}.

The example is supplied to prevent an overextension of the present
result. It is not attributed to a false theorem in the inspected
repository. In particular, one must not replace
\eqref{eq:critical-decay} by an unspecified ``at least critical decay''
in the if-and-only-if statement.

\subsection{An exponential-jet version of the test}
The residue tests can be compressed into finite exponential jets:
\begin{equation}\label{eq:jet-pencil}
 \mathcal M_{\nu,q}(t)=e^{-a_\nu t}
     \sum_{h=q}^{m_\nu}\frac{t^{h-q}}{(h-q)!}
             P_{\nu,h}(e^{-t})\pmod{t^{d+1}}.
\end{equation}
Expanding the right side shows that
\[
 \mathcal M_{\nu,q}(t)=\sum_{k=0}^d c_{\nu,q}[k]\frac{t^k}{k!}.
\]
Thus rank one means either
$\mathcal M_{\nu,q}(t)=\lambda_{\nu,q}e^{\tau t}\pmod{t^{d+1}}$
for one common $\tau$, or
$\mathcal M_{\nu,q}(t)=\lambda_{\nu,q}t^d/d!\pmod{t^{d+1}}$.
This formulation gives a short certificate representation without
constructing every matrix entry separately.
